% 本文件由 LaTeX 排版 Agent 根据有效 translation.json 自动生成。
% author=Ignatius of Antioch; work_id=0114; title=伪经：安提约基雅的圣依纳爵书信集

\chapter{伪经：安提约基雅的圣依纳爵书信集}

% structure: p0001 source=h1 target=omitted reason=page-title-already-rendered-by-parent

% structure: p0002 source=h2 target=section reason=relative-heading-rank; base=h2

\section{致塔尔索人书}

% structure: p0003 source=h3 target=subsection reason=relative-heading-rank; base=h2

\subsection{问候}

依纳爵，又称德福禄，致书给在塔尔索的教会，就是在基督内得救、堪受赞美、堪受纪念、堪受钟爱的教会：愿慈爱与平安由天主圣父和主耶稣基督，常与你们同在。\par

% structure: p0005 source=h3 target=subsection reason=relative-heading-rank; base=h2

\subsection{第一章 他自身的苦难：劝勉持守}

从叙利亚直到罗马，我与野兽搏斗：不是被残暴的野兽吞噬，因为这些野兽，如你们所知，因天主的旨意曾放过达尼尔，而是与人形的野兽搏斗，在其中隐藏着无情的野兽本身，日复一日地刺伤我、击打我。但这些艰难困苦，\BibleQuote{都不能动摇我，我也不把自己的性命视为宝贵，}\BibleRef{宗 20:24} 以至于我过于爱恋它，超过爱主。因此，我准备好（面对）烈火、野兽、刀剑或十字架，只要能看见为我死去的基督、我的救主和天主。所以我，基督的囚徒，被海陆驱赶的人，劝勉你们：\BibleQuote{你们要站稳于信德，}\BibleRef{格前 16:13} 要坚定不移，\BibleQuote{因为义人必因信德而生活；}\BibleRef{哈 2:4} 要毫不动摇，因为\Quote{主使那些具有同一品格的人同住在一座房屋里。}\par

% structure: p0007 source=h3 target=subsection reason=relative-heading-rank; base=h2

\subsection{第二章 提防谬论}

我得知某些撒殚的仆役想要扰乱你们；有的声称耶稣只是形像上诞生、形像上被钉、形像上死去；有的说他不是造物主之子；有的说他本身就是万有之上的天主。另外，有的认为他只是一个人，有的认为这肉身不会复活，所以我们应当活得放纵，享受快乐的生活，因为这正是那不久将要灭亡之人的至善。如此众多的邪恶蜂拥而至。但你们未曾\BibleQuote{顺从他们，让步一时，}\BibleRef{迦 2:5} 因为你们是保禄的同乡和门徒，他\BibleQuote{从耶路撒冷直到依里黎科，到处传布了基督的福音，}\BibleRef{罗 15:19} 并在身上带着\BibleQuote{基督的烙印。}\BibleRef{迦 6:17}\par

% structure: p0009 source=h3 target=subsection reason=relative-heading-rank; base=h2

\subsection{第三章 关于基督的真道}

怀念他，你们务必知道，主耶稣确实是出于玛利亚而诞生，由女人所生；也确实被钉在十字架上。因为他说：\BibleQuote{但我绝不以别的夸耀，只以我们主耶稣基督的十字架来夸耀。}\BibleRef{迦 6:14} 祂也确实受苦、死去并复活了。因为（保禄）说：\Quote{如果基督成为可受难的，并且第一个从死者中复活。}又说：\BibleQuote{祂死，是死于罪恶一次；祂活，是活于天主。}\BibleRef{罗 6:10} 否则，如果基督没有死，忍受束缚有什么益处？忍耐有什么益处？忍受鞭打有什么益处？为什么会有以下事实：伯多禄被钉十字架；保禄和雅各伯被刀剑杀死；若望被放逐到帕特摩斯；斯德望被那些杀害主的犹太人用石头砸死？但是（事实上）这些苦难没有一个是徒然的；因为主确实被不敬虔的人钉在十字架上。\par

% structure: p0011 source=h3 target=subsection reason=relative-heading-rank; base=h2

\subsection{第四章 续}

并且（你们还要知道），那由女人所生的，是天主子；那被钉在十字架上的，是\BibleQuote{一切受造物的首生者，}\BibleRef{哥 1:15} 也是天主圣言，祂创造了万物。因为宗徒说：\BibleQuote{只有一个天主，就是父，万物都出于祂；也只有一个主耶稣基督，万物都藉着祂而有。}\BibleRef{格前 8:6} 又说：\BibleQuote{因为只有一个天主，一个在天主与人之间的中保，就是那人基督耶稣；}\BibleRef{弟前 2:5} 并且\BibleQuote{全部是在祂内受造的：天上的和地上的，看得见的和看不见的；祂在万有之先，万有都赖祂而存在。}\BibleRef{哥 1:16}\par

% structure: p0013 source=h3 target=subsection reason=relative-heading-rank; base=h2

\subsection{第五章 驳斥前述谬误}

祂自己并非万有之上的天主和父，而是祂的子，祂（借着）说：\BibleQuote{我升到父那里，也是你们的父，升到我的天主那里，也是你们的天主。}\BibleRef{若 20:17} 又说：\BibleQuote{当万物都归属祂之后，子也要自己归属那使万物归属祂的父，好使天主成为万物之中的万有。}\BibleRef{格前 15:28} 所以，使万物归属并成为万有的是同一位，而万物所归属、并且祂自己也同万物一起归属（前者）的是另一位。\par

% structure: p0015 source=h3 target=subsection reason=relative-heading-rank; base=h2

\subsection{第六章 续}

祂也非一个纯粹的人，因为万物都是藉着祂并在他内造成的；因为\BibleQuote{万物是藉着祂造的。}\BibleRef{若 1:3} \Quote{当祂创造天时，我已在祂身旁；我同祂一起，[与祂一同]形成（世界），祂天天喜悦我。} 一个纯粹的人怎能被这样称呼：\Quote{坐在我右边？} 再者，这样的人怎能宣告：\BibleQuote{在亚巴郎未存在以前，我就存在？}\BibleRef{若 8:58} 并且，\BibleQuote{父啊，现在请以我在世界未有以前与你同享的荣耀来荣耀我。}\BibleRef{若 17:5} 有哪个人能说：\BibleQuote{我从天降下，不是为执行我的旨意，而是为执行那派我者的旨意？}\BibleRef{若 6:38} 对于哪个人，可以说：\Quote{祂是那真光，照亮一切来到世上的人；祂在世界，世界也藉着祂造成，世界却不认识祂；祂来到自己的地方，自己的人却不接受祂？} 这样的人怎会是纯粹的人，从玛利亚开始存在，而不是天主圣言和独生子呢？因为\BibleQuote{在起初已有圣言，圣言与天主同在，圣言就是天主。}\BibleRef{若 1:1} 在别处又说：\Quote{主在我造化之先创造了我，为祂的工程，在太古之初，在群山尚未被创造以前，我已受生。}\par

% structure: p0017 source=h3 target=subsection reason=relative-heading-rank; base=h2

\subsection{第七章 续}

祂也显示我们的身体将要复活，因为祂说：\Quote{我实在告诉你们，时候将到，那时凡在坟墓里的都要听见天主圣子的声音；听见的人就必活着。} \newline{} 宗徒也说：\Quote{因为这必朽坏的必须穿上不朽坏的，这必死的必须穿上不死的。} \BibleRef{格前 15:53} \newline{} 我们应当度清醒和公义的生活，祂又这样说：\Quote{不要自欺：无论是奸淫者、作娈童的、与男人苟合的、行淫的、辱骂人的、醉酒的、偷窃的，都不能继承天主的国。} \BibleRef{格前 6:9} \newline{} 又说：\Quote{如果死人不复活，基督也就没有复活；那么，我们的宣讲便是空的，你们的信仰也是空的；你们仍在你们的罪中。那么，在基督内睡了的人也丧失亡了。如果我们只在今生寄望于基督，我们就是众人中最可怜的了。如果死人不复活，我们就吃吃喝喝吧，因为明天就要死了。} \BibleRef{格前 15:13} \newline{} 但如果我们的状况和心意是这样，我们与驴和狗有什么分别？它们不关心将来，只想到吃，并满足吃后随之而来的情欲，因为在他们内没有任何理智运行。\par

% structure: p0019 source=h3 target=subsection reason=relative-heading-rank; base=h2

\subsection{第8章  劝勉圣善与秩序}

愿我在主内因你们而喜乐！你们要清醒。你们每个人都要丢弃一切恶意、兽般的愤怒、诽谤、诬蔑、污言、秽语、耳语、傲慢、醉酒、淫欲、贪婪、虚荣、嫉妒，以及一切类似的行为。\Quote{但要穿上主耶稣基督，不要为肉体的私欲作准备。} \BibleRef{罗 13:14} \newline{} 你们诸位司铎，要服从主教；你们诸位执事，要服从司铎；而你们众人，要服从司铎和执事。愿我的灵魂为那些保存这良好秩序的人而存在；愿主常与他们同在！\par

% structure: p0021 source=h3 target=subsection reason=relative-heading-rank; base=h2

\subsection{第9章  劝勉履行相关职责}

你们做丈夫的，要爱你们的妻子；你们做妻子的，要爱你们的丈夫。你们做子女的，要孝敬父母。你们做父母的，\Quote{要以主的管教和训诫养育你们的孩子。} \BibleRef{弗 6:4} \newline{} 要尊敬那些持守童贞的人，如同基督的祭司；也要尊敬那些在庄重品行上持守的寡妇，如同天主的祭台。你们做仆人的，要以敬畏的心服侍主人。你们做主人的，要以温良向仆人发号施令。你们中不可有人闲懒，因为闲懒是匮乏之母。我吩咐这些事，并非因为我是什么重要人物，虽然我为基督而被捆锁；而是作为兄弟，我提醒你们这些事。愿主与你们同在！\par

% structure: p0023 source=h3 target=subsection reason=relative-heading-rank; base=h2

\subsection{第10章  问候}

愿我能享有你们的祈祷！请为我祈祷，使我得以到达耶稣那里。我把在安提约基雅的教会交托给你们。斐理伯的教会——我也从那里给你们写信——问候你们。你们的执事斐洛，我因他热心地在我一切事上服侍而感谢他，问候你们。从叙利亚来的执事阿加托珀斯，他在基督内跟随我，问候你们。\Quote{你们要以圣洁的吻彼此问候。} \BibleRef{伯前 5:14} \newline{} 我问候你们所有在基督内的男女。愿你们在身体、灵魂和同一圣神内平安，不要忘记我。愿主与你们同在！\par

% structure: p0025 source=h2 target=section reason=relative-heading-rank; base=h2

\section{致安提约基雅人书}

% structure: p0026 source=h3 target=subsection reason=relative-heading-rank; base=h2

\subsection{问候}

依纳爵，又名塞奥弗鲁斯，致书给在叙利亚作客的教会，这教会从天主获得了怜悯，被基督拣选，并且首先\BibleRef{宗 11:26}被称为基督的名。愿你们在父天主和主耶稣基督内获得幸福。\par

% structure: p0028 source=h3 target=subsection reason=relative-heading-rank; base=h2

\subsection{第1章  谨防谬误}

主使我的锁链变得轻松容易，自从我知道你们在平安中，在肉体和精神上都生活在完全和谐里。\Quote{因此我这为主被囚的劝你们，行事为人要与你们所蒙的召叫相称，} \BibleRef{弗 4:1} \newline{} 防备那些侵入我们中间、欺骗并毁灭接受它们之人的魔鬼的异端；你们要注意宗徒的教训，并相信法律和先知；要拒绝一切犹太人和外邦人的谬误，既不要引入多神，也不要借口保持天主的唯一性而否认基督。\par

% structure: p0030 source=h3 target=subsection reason=relative-heading-rank; base=h2

\subsection{第2章  关于天主和基督的真道}

因为天主的忠仆梅瑟，当他说：\Quote{上主你的天主是唯一的上主，} \BibleRef{申 6:4} 这样宣告只有一位天主时，却也立刻承认了我们的主，他说：\Quote{上主从天上上主那里降下火和硫磺，降在所多玛和蛾摩拉上。} \BibleRef{创 19:24} \newline{} 又说：\Quote{天主说：让我们按我们的肖像造人；天主就造了人，按天主的肖像造了他。} \BibleRef{创 1:26} \newline{} 又说：\Quote{按天主的肖像造了人。} \BibleRef{创 5:1} \newline{} 至于[天主子]要成为人，[梅瑟]说：\Quote{上主要从你们弟兄中给你们兴起一位像我一样的先知。}\par

% structure: p0032 source=h3 target=subsection reason=relative-heading-rank; base=h2

\subsection{第3章  接续上文}

先知们以天主的口气说话：\Quote{我是天主，我是首先的，也是末后的，除我以外没有别的天主，} \BibleRef{依 44:6} 这是论及万物之父，但也论及我们的主耶稣基督。他们说：\Quote{有一个儿子赐给了我们，他的肩上有统治权；他的名字被称为大谋略的天使、奇妙、策士、大能而有力的天主。} \BibleRef{依 9:6} \newline{} 关于祂的降生成人：\Quote{看，一位童贞女将怀孕生子，人要称祂的名字为厄玛奴耳。} \BibleRef{依 7:14} \newline{} 关于祂的苦难：\Quote{祂被牵去如羊待宰，又如羔羊在剪毛者前无声，我也像一只无辜的羔羊被牵去祭杀。} \BibleRef{依 53:7}\par

% structure: p0034 source=h3 target=subsection reason=relative-heading-rank; base=h2

\subsection{第4章  接续上文}

福音作者们宣告圣父是\Quote{唯一真天主}\BibleRef{若 17:3}时，也没有忽略关于主的事，反而写道：\Quote{在起初已有圣言，圣言与天主同在，圣言就是天主。这圣言起初与天主同在。万物是藉着祂造成的；凡受造的，没有一样不是由祂造成的。}\BibleRef{若 1:1}关于降生成人：[圣经]说：\Quote{圣言成了血肉，住在我们中间。}\BibleRef{若 1:14}又说：\Quote{亚巴郎之子，达味之子耶稣基督的族谱。}\BibleRef{玛 1:1}那些曾说过\Quote{天主只有一个}\BibleRef{格前 8:4}的宗徒们，也说了\Quote{在天主与人之间只有一个中保}\BibleRef{弗 4:5}。他们并不以降生和受难为耻。因为[一位]说什么？\Quote{基督耶稣，祂把自己交出}\BibleRef{弟前 2:5}为世界的生命和救恩。\par

% structure: p0036 source=h3 target=subsection reason=relative-heading-rank; base=h2

\subsection{第五章  斥责假教师}

因此，凡宣告只有一位天主，却以此否认基督天主性的人，就是魔鬼，是一切正义的仇敌。凡承认基督，却不承认祂是造物主之子的，而是某个律法和先知所宣告者之外的未知存有之子，这人就是魔鬼的工具。凡否认降生成人，并以十字架为耻（我为这十字架而受束缚）的人，就是假基督。此外，凡肯定基督仅是凡人的人，按先知的话\BibleRef{耶 17:5}是受诅咒的，因为他不信靠天主，而信靠人。因此，他也像野生的香桃木一样不结果实。\par

% structure: p0038 source=h3 target=subsection reason=relative-heading-rank; base=h2

\subsection{第六章  再次警告}

我写这些话给你们，基督的新橄榄树，并非我知晓你们持有任何此类意见，而是为提醒你们，如同父亲提醒他的儿女。因此，要提防那些急于行恶的人，那些\Quote{基督十字架的仇敌，他们的结局是丧亡，他们的光荣在于他们的羞耻}\BibleRef{斐 3:18}。要提防那些\Quote{哑巴狗}，那些拖曳的蛇，那些有鳞的龙，那些毒蛇、蜥蜴和蝎子。因为它们是狡猾的狼，是模仿人形的猿猴。\par

% structure: p0040 source=h3 target=subsection reason=relative-heading-rank; base=h2

\subsection{第七章  劝勉行为一致}

你们曾是保禄和伯多禄的门徒；不要失去交托给你们的。要纪念欧奥迪亚，你们理应受祝福的牧者，宗徒们首次将治理你们的职责托付给他。我们不要让我们的父蒙羞。让我们证明自己是祂的真子女，而非私生子。你们知道我在你们中间是怎样行事的。我在时对你们所说的话，我在远方时现在也写信给你们。\Quote{若有人不爱主耶稣基督，他就该受咒诅。}\BibleRef{格前 16:22}你们要效法我。\BibleRef{格前 4:16}当我到达耶稣那里时，愿我的灵魂代替你们的。请记念我的锁链。\BibleRef{哥 4:18}\par

% structure: p0042 source=h3 target=subsection reason=relative-heading-rank; base=h2

\subsection{第八章  对长老和其他人的劝勉}

你们这些长老，\Quote{牧放你们中间的天主的羊群}\BibleRef{伯前 5:2}，直到天主指示谁将管理你们。因为\Quote{我已被奠祭}\BibleRef{弟后 4:6}，好使我\Quote{赢得基督}\BibleRef{斐 3:8}。让执事们知道他们的尊严，并努力无可指摘，成为基督的追随者。让子民顺从长老和执事。让贞女们知道她们已将自己奉献给谁。\par

% structure: p0044 source=h3 target=subsection reason=relative-heading-rank; base=h2

\subsection{第九章  丈夫、妻子、父母和子女的本分}

让丈夫爱妻子，记念在创造时，一个男人只被赐予一个女人，而非多个。让妻子尊敬丈夫，如同自己的肉身；不要让她们敢于直呼丈夫的名字。\BibleRef{伯前 3:6}让她们保持贞洁，视丈夫为唯一的伴侣，按天主的旨意与丈夫结合。你们作父母的，要以神圣的教训教养你们的子女。你们作子女的，\Quote{要孝敬父母，使你得福。}\par

% structure: p0046 source=h3 target=subsection reason=relative-heading-rank; base=h2

\subsection{第十章  主人和仆人的本分}

你们作主人的，不要傲慢地对待仆人，而要效法忍耐的约伯，他宣告说：\Quote{我岂轻视我的仆人或婢女，当他们与我争论时？否则，上主追究我时，我该怎么办？}\BibleRef{约 31:13}你们也知道接下来的话。你们作仆人的，不要在凡事上激怒主人，免得你们自招不治之祸。\par

% structure: p0048 source=h3 target=subsection reason=relative-heading-rank; base=h2

\subsection{第十一章  各种道德责任的训导}

凡懒于作工的人不可吃饭，\BibleRef{得后 3:10}免得他闲游浪荡，成为奸淫者。醉酒、愤怒、嫉妒、辱骂、争吵、亵渎，\Quote{连提也不可提}\BibleRef{弗 5:3}。不要让寡妇过享乐的生活，免得她们悖逆圣经。\BibleRef{弟前 5:6}在一切不涉及[属灵]危险的事上，要顺服凯撒。不要激怒管理你们的人，免得给那寻找机会的人留下把柄。至于行邪术、男色或杀人，写信给你们已是多余，因为这些恶行连外邦人也禁止。我并非作为宗徒就这些事发号施令，而是作为你们的同仆提醒你们。\par

% structure: p0050 source=h3 target=subsection reason=relative-heading-rank; base=h2

\subsection{第十二章  问安}

我问候圣洁的长老团。我问候神圣的执事们，以及那位我最亲爱的人，愿我藉着圣神看见他代替我的职位，当我到达基督那里时。愿我的灵魂代替他的。我问候副执事、读经者、唱诗者、守门者、劳作者、驱魔者、宣信者。我问候圣洁的门卫、在基督内的女执事。我问候许配给基督的贞女，愿我在主耶稣里因她们而喜乐。我问候主的子民，从最小的到最大的，以及主内所有的姊妹。\par

% structure: p0052 source=h3 target=subsection reason=relative-heading-rank; base=h2

\subsection{第十三章  问安继续}

我问候卡西安（Cassian）及其生活伴侣，以及他们亲爱的子女。最可敬的主教波利卡普（Polycarp）也深切关心你们，他向你们致候；我在主内将你们托付给了他。士每拿整个教会确实在祈祷中于主内纪念你们。厄弗所的牧者敖乃息摩（Onesimus）向你们致候。玛格内西的主教达玛斯（Damas）向你们致候。特拉里亚的主教波里比乌（Polybius）向你们致候。我的同伴，执事斐罗（Philo）和阿加托普斯（Agathopus）向你们致候。\Quote{你们要以圣吻彼此问候。}\BibleRef{格后 13:12}\par

% structure: p0054 source=h3 target=subsection reason=relative-heading-rank; base=h2

\subsection{第十四章 结论}

我从斐理伯写这封信给你们。愿那\OriginalTerm{无始者}{unbegotten}藉着那在时间开始前\OriginalTerm{受生者}{begotten}使你们在灵魂和肉体上坚定不移！愿我在基督的王国中见到你们！我问候那将代替我管辖你们的人：愿我在主内因他而喜乐！你们要藉着圣神的光照，在天主内并基督内安好。\par

% structure: p0056 source=h2 target=section reason=relative-heading-rank; base=h2

\section{致赫罗书——安提约基雅的执事}

% structure: p0057 source=h3 target=subsection reason=relative-heading-rank; base=h2

\subsection{问候}

依纳爵，亦称为德奥福禄，致基督的执事、天主的仆人赫罗，一位蒙天主尊敬、深受爱戴与敬重的人，他怀有基督和圣神，是我在信德和爱德中的儿子：愿恩宠、仁慈与平安从全能的天主，并从我们的主耶稣基督、祂的独生子那里赐予你，\Quote{祂为我们舍弃了自己，为救赎我们脱离这邪恶的现世，}\BibleRef{迦 1:4} 并保守我们进入祂的天国。\par

% structure: p0059 source=h3 target=subsection reason=relative-heading-rank; base=h2

\subsection{第一章 劝勉热诚与节制}

我在天主内劝勉你，要加速你的进程，并维护你的尊严。留意与圣徒保持和谐。要担当软弱者的负担，好使\Quote{你满全基督的法律。}\BibleRef{迦 6:2} 致力于禁食与祈祷，但不可过度，免得因此毁了自己。不要完全戒酒戒肉，因为这些事物并非可憎的，因为（圣经）说：\Quote{你们要吃地上的美物。}\BibleRef{依 1:19} 又说：\Quote{你们要像青菜一样吃肉。}\BibleRef{创 9:3} 又说：\Quote{酒使人心快乐，油使人容光焕发，粮食使人健壮。} 但一切都要节制使用，因为都是天主的恩赐。\Quote{因为若非出于祂，谁能吃，谁能喝？因为凡是美丽的，都是属于祂的；凡是美好的，都是属于祂的。} 要注重诵读，\BibleRef{弟前 4:13} 不仅使你本人认识律法，也能向他人解释，做天主的忠实仆人。\Quote{凡当兵的，不将世务缠身，为要使那招他当兵的人喜悦；若有人竞赛，除非按规矩竞赛，也不能得冠冕。}\BibleRef{弟后 2:4} 我这带锁链的人祈祷，愿我的灵魂代替你们的位置。\par

% structure: p0061 source=h3 target=subsection reason=relative-heading-rank; base=h2

\subsection{第二章 警惕假教师}

凡教导超乎所命令的，即使他被视为可信，即使他惯常禁食，即使他生活节制，即使他行奇迹，即使他有说预言的恩赐，你应视他如披着羊皮的狼，\BibleRef{玛 7:15} 致力于灭亡羊群。若有人否认十字架，以苦难为耻，你应视他如仇敌本身。\Quote{即使他把自己所有的财产分施给穷人，即使他移山，即使他把身体焚烧，}\BibleRef{格前 13:2} 你应视他为可憎的。若有人轻视基督来临时所完成的律法或先知，你应视他为\OriginalTerm{假基督}{antichrist}。若有人说主只是一个凡人，他就是犹太人，是基督的谋杀者。\par

% structure: p0063 source=h3 target=subsection reason=relative-heading-rank; base=h2

\subsection{第三章 关于教会职务的劝勉}

\Quote{要尊敬那真正是寡妇的寡妇。}\BibleRef{弟前 5:3} 要做孤儿的友人，因为天主是\Quote{孤儿的父亲，寡妇的审判者。}不可不经主教而行任何事，因为他们是司铎，你是司铎的仆人。他们施洗、献祭、祝圣并覆手；你则服务于他们，如同圣斯德望在耶路撒冷服务于雅各伯和长老们一样。不可忽略神圣的集会；要一一询问每个人的名字。\Quote{不要叫人小看你年轻，却要在言语和行为上作信徒的榜样。}\BibleRef{弟前 4:12}\par

% structure: p0065 source=h3 target=subsection reason=relative-heading-rank; base=h2

\subsection{第四章 不可轻视仆人与妇女}

不要以仆人为羞耻，因为我们与他们有同样的本性。不要憎恨妇女，因为她们生养了你。因此，应当爱那些生养我们的人（但只能在主内），因为人没有女人就不能生育子女。所以，我们应当尊敬那些参与生育我们的人。\Quote{男人没有女人，女人没有男人，也不可分开，}\BibleRef{格前 11:11} 除了最初被造的那一对。因为亚当的身体由\OriginalTerm{四元素}{four elements}形成，厄娃的身体出自亚当的肋骨。而且，主的出生极为独特，仅由一位童贞女所生。这并非因为合法结合（男女）是可憎的，而是这样的出生方式适合于天主。因为\OriginalTerm{造物主}{Creator}不宜采用通常的生育方法，而应采用独特且奇妙的方法，因为祂是造物主。\par

% structure: p0067 source=h3 target=subsection reason=relative-heading-rank; base=h2

\subsection{第五章 各种相对本分}

要逃避傲慢，\Quote{因为天主抵挡骄傲的人。}\BibleRef{雅 4:6} 要憎恶虚伪，因为（圣经）说：\Quote{你要消灭一切说谎的人。}要警惕嫉妒，因为嫉妒的创始者是魔鬼，而魔鬼的继承者是加音，他嫉妒自己的兄弟，因嫉妒而杀人。要劝勉我的姐妹们爱天主，并只满足于自己的丈夫。同样，也要劝勉我的兄弟们只满足于自己的妻子。要照看贞女，如同基督的宝贵珍宝。要宽容忍耐，\BibleRef{箴 14:29} 使你在智慧上伟大。不可忽略穷人，只要你能做到。因为\Quote{施舍和忠诚能赎罪。}\par

% structure: p0069 source=h3 target=subsection reason=relative-heading-rank; base=h2

\subsection{第六章 劝勉纯洁与谨慎}

你要保持洁净，如同天主的居所。你是基督的圣殿。你是圣神的器皿。你知道我怎样培育了你。虽然我是最卑微的人，仍要努力效法我，作我行为的模仿者。我不在世上夸耀，只夸耀主。我劝勉我的儿子\Person{赫罗}：\BibleQuote{但那夸耀的，应当夸耀主。} \BibleRef{格前 1:31} 愿我因你而喜乐，我亲爱的儿子，愿那唯一无生的天主和主耶稣基督保守你！不要相信所有的人，不要信赖所有的人；也不要让任何人以谄媚胜过你。因为有许多是撒殚的仆役，而\BibleQuote{轻信的人，心是轻浮的。} \BibleRef{德 19:4}\par

% structure: p0071 source=h3 target=subsection reason=relative-heading-rank; base=h2

\subsection{第七章：对\Person{赫罗}的庄严托付，作为未来\Place{安提约基雅}的主教}

要常记念天主，你就永不犯罪。祈祷时不要心怀二意\BibleRef{雅 1:6}，因为那不疑虑的人是有福的。因为我相信主耶稣基督的父和祂的独生子，天主必会在我（的宝座上）显示\Person{赫罗}。所以，你要加速你的进程。我在宇宙的天主面前，并在基督面前，在圣神和服役的（天使）层级面前，嘱咐你：要安全地保管我和基督所托付给你的那\OriginalTerm{寄托}{deposit}，不要因那些天主（向我）所显示关于你的事而自认不配。我把\Place{安提约基雅}的教会交给你。我已在主耶稣基督里把你托付给\Person{波利卡普}。\par

% structure: p0073 source=h3 target=subsection reason=relative-heading-rank; base=h2

\subsection{第八章：问安}

主教们——\Person{奥内西慕斯}、\Person{比图斯}、\Person{达玛斯}、\Person{波利比乌斯}，以及所有在\Place{斐理伯}的（我也从那里写信给你们）——在基督里问候你。请问候那堪当天主的\OriginalTerm{长老团}{presbytery}；请问候我神圣的同工执事们，愿我在基督里因他们而喜乐，无论是肉身还是心灵。请问候主的子民，从最小的到最大的，每个都提名；我把他们托付给你，如同\Person{梅瑟}（将以色列人）托付给继他之后的领袖\Person{若苏厄}。不要以为我这话是狂妄；因为我们虽不如他们，但至少我们祈祷能如此，因为我们确是亚巴郎的子孙。所以，\Person{赫罗}，要刚强如英雄，如男子汉。因为从今以后，你要\BibleRef{申 31:7}带领\Place{安提约基雅}主的子民出入，这样，\BibleQuote{主的会众就不致像没有牧人的羊群。} \BibleRef{户 27:17}\par

% structure: p0075 source=h3 target=subsection reason=relative-heading-rank; base=h2

\subsection{第九章：结束的问候与指示}

请问候我的主人\Person{卡西安}，和他极其稳重的伴侣，以及他们非常可爱的孩子们；愿\BibleQuote{主使他们在那一日从主那里获得怜悯，} \BibleRef{弟后 1:18}因他们对我们的服侍；我也在基督里将他们托付给你。请问候在\Place{劳狄刻雅}所有在基督里的信徒，各提名。不要忽略在\Place{塔尔索}的人，而要坚定地照管他们，在福音里坚固他们。我在主里问候\Place{纳波利}（靠近\Place{阿纳扎尔布斯}）的主教\Person{马里斯}。也请问候我的女儿\Person{玛利亚}，她因庄重和学识而出众，也问候\BibleQuote{在她家的教会。} \BibleRef{哥 4:15} 愿我的灵魂取代她的灵魂：她是虔诚妇女的榜样。愿基督的父，藉祂的独生子，保守你健康，在一切事上声誉卓著，直到极老之年，为了天主的教会的益处！愿你在主里安好，并祈祷使我得以成全。\par

% structure: p0077 source=h2 target=section reason=relative-heading-rank; base=h2

\section{致斐理伯人书}

% structure: p0078 source=h3 target=subsection reason=relative-heading-rank; base=h2

\subsection{问候}

\Person{伊纳爵}，又称\Person{德福禄}，致书给在\Place{斐理伯}的天主的教会，这教会藉信德、忍耐和\OriginalTerm{无伪的爱}{love unfeigned}获得了怜悯：愿从天主父和主耶稣基督——\BibleQuote{祂是众人的救主，特别是信徒的救主}——而来的怜悯与平安归于你们。 \BibleRef{弟前 4:10}\par

% structure: p0080 source=h3 target=subsection reason=relative-heading-rank; base=h2

\subsection{第一章：写作此书的缘由}

想起你们在基督里的爱和热忱（你们曾向我们显明），我们认为适宜写信给你们，你们对弟兄们显示如此敬虔和属灵的爱，为要提醒你们的基督徒路程，\Quote{使你们众人都说一样的话，同心合意，怀有同样的思想，并照同一信德的规则行事，}如同\Person{保禄}所劝诫的。因为若只有一位宇宙的天主、基督的父，\BibleQuote{万物都出于祂；} \BibleRef{格前 8:6}并只有一位主耶稣基督、我们的[主]，\BibleQuote{万物都藉着祂；} \BibleRef{格前 8:6}也只有一位圣神，祂曾在\Person{梅瑟}、先知和宗徒内运行\BibleRef{格前 12:11}；也只有一次圣洗，它被施行是为使我们与主的死亡有分；也只有一位\OriginalTerm{蒙选的教会}{elect Church}；同样，对基督也应该只有一种信德。因为\BibleQuote{只有一个主，一个信德，一个圣洗；一个天主，就是众人的父，祂超越众人，贯穿众人，且在众人之内。} \BibleRef{弗 4:5}\par

% structure: p0082 source=h3 target=subsection reason=relative-heading-rank; base=h2

\subsection{第二章：天主圣三的合一}

所以，只有一位天主、父，不是两位或三位；那位自有者，除祂以外没有别的——唯一真实的天主。因为经上记载：\BibleQuote{上主你的天主是唯一的主。}\BibleRef{申 6:4} 又说：\BibleQuote{不是只有一位天主创造了我们？我们岂不都同有一位父亲？}\BibleRef{拉 2:10} 也只有一位圣子——天主圣言。因为经上说：\BibleQuote{独生子，在父怀里的。}\BibleRef{若 1:18} 又说：\BibleQuote{只有一位主耶稣基督。}\BibleRef{格前 8:6} 另一处说：\BibleQuote{祂叫什么名字？祂的儿子叫什么名字，你可知？}\BibleRef{箴 30:4} 也只有一位护慰者。因为经上说：\BibleQuote{只有一个圣神。}\BibleRef{弗 4:4} 既然\BibleQuote{我们蒙召，同有一个希望。}\BibleRef{格前 12:13} 又说：\BibleQuote{我们都因一个圣神而饮。}\BibleRef{弗 4:4} 以及接下来的话。显而易见，所有这些〔信徒所领受的〕恩赐，\BibleQuote{都是同一圣神所运行。}\BibleRef{格前 12:11} 所以不是三位父、三位子或三位护慰者，而是一位父、一位子和一位护慰者。因此，主派遣宗徒去使万民成为门徒时，吩咐他们\BibleQuote{因父及子及圣神之名施洗}，\BibleRef{玛 28:19} 不是归入一个有三个名字的位格，也不是归入三个成了血肉的位格，而是归入三个同样尊荣的位格。\par

% structure: p0084 source=h3 target=subsection reason=relative-heading-rank; base=h2

\subsection{基督真实降生，并受死}

因为只有一位成了血肉，既不是父也不是护慰者，唯独是圣子——并非外表或想象，而是真实地成了血肉。因为\BibleQuote{圣言成了血肉。}\BibleRef{若 1:14} 因为\BibleQuote{智慧建造了房屋。}\BibleRef{箴 9:1} 天主圣言由童贞女取了肉身降生为人，没有男子的交合。因为经上记载：\BibleQuote{童贞女要怀孕生子。}\BibleRef{依 7:14} 祂确实降生，确实成长，确实吃、喝，确实被钉、死亡，并复活。信这些事真实发生、真实成就的人是有福的。不信的人与钉主的人同受诅咒。因为今世的首领见人否认十字架便欢喜，因他知道承认十字架就是他的毁灭。因为那是竖立起来反对他权势的凯旋标记，他一看见就战栗，一听到便惧怕。\par

% structure: p0086 source=h3 target=subsection reason=relative-heading-rank; base=h2

\subsection{撒殚的恶毒与愚蠢}

确实，在十字架竖立之前，他（撒殚）切望此事；他\BibleQuote{在悖逆之子身上运行}\BibleRef{弗 2:2}，运行在犹达斯、法利塞人、撒杜塞人、老年人、少年人及司祭身上。但当十字架即将竖立时，他惊慌了，便给叛徒注入悔意，指给他上吊的绳索，教他缢死。他又惊吓那愚蠢的妇人，以梦境扰乱她；那曾用尽一切手段预备十字架者，如今却试图阻止其竖立；并非因为他罪大恶极而悔改（若如此他便不会全然败坏），而是因他察觉自己的毁灭就在眼前。因为基督的十字架是他定罪的开端、死亡的开端、毁灭的开端。故此，他在某些人心中运行，使他们否认十字架、羞耻于苦难、称死亡为幻象、篡改并曲解童贞女的生育、并毁谤人性本身为可憎。他与犹太人一同作战否认十字架，与外邦人一同作战毁谤玛利亚——那些持异端主张基督仅有幻影身体的人。因为一切邪恶的首领以多形出现，他是人类的欺骗者，反复无常，甚至自相矛盾，计划一事后又行另一事。他智慧于作恶，却全然不知何为善。确实，他满是无知，因他故意缺乏理性：一个人怎能被视为别的，当他连脚边的理性都辨不出？\par

% structure: p0088 source=h3 target=subsection reason=relative-heading-rank; base=h2

\subsection{对撒殚的呼告}

因为若主只是个有灵魂和身体的人，你为何篡改并曲解祂按人性共通本性而生？你为何称苦难为幻象，仿佛发生在〔一个〕人身上是奇怪的事？你为何将必死者的死亡仅视为虚幻的死亡？但若祂既是天主又是人，那么为何你称祂为\BibleQuote{荣耀的主}\BibleRef{格前 2:8}（祂本性不变）是不合法的？为何你说宣告那拥有灵魂的立法者\BibleQuote{圣言成了血肉}\BibleRef{若 1:14}、且成为完全的人、而不仅仅是居住在人里面的那位，是不合法的？但这行法术者如何存在？祂从太古按父的旨意形成一切感知和理智可把握的自然，当祂成为血肉时，治愈了各种疾病和软弱。\BibleRef{玛 4:23}\par

% structure: p0090 source=h3 target=subsection reason=relative-heading-rank; base=h2

\subsection{续}

祂岂不真是天主吗？祂使死人复活，使瘸子四肢健全，洁净癞病人，使瞎子复明，增加或改变已有的物质，如同五个饼和两条鱼，以及变成酒的水，并仅凭一句话就使你的全军溃败。你为什么辱骂\OriginalTerm{童贞女}{Virgin}的本性，称她的肢体可耻，既然你从前在公开游行中展示它们，命令它们赤裸展出，男性在女性面前，女性激起男性无节制的情欲？但现在你却认为这些是可耻的，你假装充满谦逊，你这淫乱之神，不知道只有被邪恶玷污时，事物才变得可耻。但当罪不存在时，被造之物无一可耻，无一邪恶，全都甚好。但因你盲目，你辱骂这些事物。\par

% structure: p0092 source=h3 target=subsection reason=relative-heading-rank; base=h2

\subsection{第七章 继续：撒殚的自相矛盾}

再者，基督在你看来如何完全不是由童贞女所生，而是万有之上的天主、全能者？那么请说，谁派遣了祂？谁是祂的主宰？祂服从谁的旨意？祂履行了什么法律？既然祂并不服从任何人的意志或权能。当你否认基督降生时，你肯定那未受生者受生，那无始者被钉在十字架上，我不知道是经谁的许可。但你的多变策略瞒不过我，我也知道你会走歪斜不定的步子。你自夸知晓一切，却不知道真正降生的是谁。\par

% structure: p0094 source=h3 target=subsection reason=relative-heading-rank; base=h2

\subsection{第八章 继续：撒殚的无知}

有许多事是你不知道的：[如下]：玛利亚的\OriginalTerm{童贞}{virginity}；奇妙的诞生；\OriginalTerm{降生成人}{Incarnation}的是谁；引导东方贤士的星；奉献礼物的\OriginalTerm{贤士}{Magi}；总领天使对童贞女的问候；那已订婚者的奇妙受孕；那孩童前驱关于童贞女子所发的宣告，以及他因预见到的事而在母腹中跳跃；天使们为那降生者歌唱；报告给牧羊人的喜讯；黑落德怕被夺去王位的恐惧；屠杀婴孩的命令；迁往埃及，并从那里回到同一地区；婴孩的襁褓；人口登记；以母乳喂养；称那未生育者为父亲；由于无处容身而用马槽；（为婴孩）没有人的预备；逐渐成长、人的语言、饥饿、口渴、旅途、疲惫；献祭，然后还有割损、受洗；天主在受洗者之上的声音，论到祂是谁，从何处来；圣神和圣父从上而来的见证；先知若翰的声音，当祂以\Quote{羔羊}之称预示\OriginalTerm{苦难}{Passion}时；行各种奇迹，诸多医治；主斥责掌管海洋和风；驱逐邪神；你自己受折磨，当你被你显现的那位的能力打击时，你无能为力。\par

% structure: p0096 source=h3 target=subsection reason=relative-heading-rank; base=h2

\subsection{第九章 继续：撒殚的无知}

看到这些事，你完全困惑。你不知道应是童贞女生产；但天使的赞美歌令你惊讶，贤士的朝拜和星的显现也是如此。你回到（故意）无知的状态，因为你认为这一切都微不足道；你以为襁褓、割损、母乳喂养是可鄙的：这些事在你看来不配天主。再者，你看见一个人四十昼夜不尝人间食物，有服役的天使同在，他们的出现令你战栗，当初你先看见祂像普通人一样受洗，你不知道原因。但在祂（长久）禁食之后，你又恢复惯常的胆大，在祂饥饿时试探祂，好像祂是普通人，不知道祂是谁。因为你说：\Quote{你若是天主子，就命这些石头变成饼吧。}\BibleRef{玛 4:3} 现在，\Quote{你若是子}这个说法，是无知的标志。因为你若真有知识，就会明白造物主能同样容易地创造不存在的，并改变已有的事物。你借着饥饿试探那供养一切需要食物者。你试探了那\Quote{光荣的主}\BibleRef{格前 2:8}，在你的恶意中忘记了\Quote{人生活不只靠饼，也靠天主口中所出的一切言语}。因为若你知道祂是天主子，你也会明白那使身体四十昼夜不觉缺乏的，也能永远如此。那么，祂为何忍受饥饿？为证明祂取了与普通人同样感受的身体。由第一件事祂显示祂是天主，由第二件显示祂也是人。\par

% structure: p0098 source=h3 target=subsection reason=relative-heading-rank; base=h2

\subsection{第十章 继续：撒殚的胆大}

那么，你这从最高荣耀\Quote{如闪电}\BibleRef{路 10:18}坠落的，竟敢对主说：\Quote{你从这里跳下去}\BibleRef{玛 4:6}，对那使无变为有的\BibleRef{罗 4:17}，并激怒那毫无虚荣的祂炫耀虚妄的荣耀吗？你岂不假装在圣经中读到关于祂的话：\Quote{祂为你吩咐了祂的天使，他们要用手托着你，免得你的脚碰在石头上。}\BibleRef{玛 4:6} 同时你假装不知道其余的话，偷偷隐藏了（圣经）关于你和你的仆人所预言的：\Quote{你要践踏蝮蛇和毒蜥，你要踩踏狮子和毒龙。}\par

% structure: p0100 source=h3 target=subsection reason=relative-heading-rank; base=h2

\subsection{第十一章 继续：撒殚的胆大}

因此，倘若你们被践踏在主的脚下，你们怎敢试探那不可被试探者，忘记了立法者的那条诫命：\BibleQuote{你们不可试探主你们的天主？}\BibleRef{申 6:16}是啊，你这最可咒骂的，甚至胆敢将天主的工程据为己有，宣称对这些事物的统治权已交付给你。\BibleRef{路 4:6}你还将自己的堕落作为榜样摆在主面前，应许祂只要俯伏朝拜你，就将本来属于祂的赐给祂。\BibleRef{玛 4:9}你这因恶意而比一切邪灵更邪恶的灵，怎不战栗，竟对主说出这样的话？你因欲望而跌倒，因虚荣而蒙羞；因贪婪和野心，你如今将他人拖入不敬神。你这贝里雅耳，巨龙，叛徒，弯曲的蛇，造反的天主者，被基督驱逐，与圣神隔绝，被天使行列放逐，诋毁天主的法律，一切合法之物的仇敌，你曾起来攻击最初受造的人，将那从未伤害过你的人赶出诫命；你曾兴起杀害亚伯的该隐；你曾拿起武器攻击约伯：你对主说：\BibleQuote{你若俯伏朝拜我？}何等胆大！何等疯狂！你这逃跑的奴隶，无可救药的奴隶，你竟反抗良善的主？你对如此伟大的主，那一切可感知与可理解之物的天主，说：\BibleQuote{你若俯伏朝拜我？}\par

% structure: p0102 source=h3 target=subsection reason=relative-heading-rank; base=h2

\subsection{第十二章　基督温良的答复}

但主宽容忍耐，并未将无知而胆敢说出这种话的人消灭，而是温良地回答说：\BibleQuote{去吧，撒殚。}\BibleRef{玛 4:10}祂没有说：\Quote{退到我后面去，}因为他不可能悔改；而是说：\BibleQuote{去吧，撒殚，}到你选择的那条路上去。\BibleQuote{去吧，}回到那些因你的恶意所招致的事物中去。因为我知道我是谁，是谁派遣了我，我应当朝拜谁。因为\BibleQuote{你要朝拜主你的天主，惟独事奉祂。}\BibleRef{玛 4:10}我认识那独一者；我知晓那独一的主，是你从他那里成了叛徒。我不是天主的仇敌；我承认他的卓越；我认识那是我生命之源的父。\par

% structure: p0104 source=h3 target=subsection reason=relative-heading-rank; base=h2

\subsection{第十三章　各种劝勉与指示}

弟兄们，出于我对你们的关爱，我感到必须写这些，为天主的荣耀而劝勉你们，并非因为我是什么重要人物，而只是作为一个弟兄。要服从主教、司铎和执事。要在主内彼此相爱，如同天主的肖像。你们做丈夫的，要爱你们的妻子如同自己的肢体。你们做妻子的，要爱你们的丈夫，因你们结合而与他们成为一体。若有人守贞洁或节制，不要因此自高，免得失去赏报。不要轻看庆节。不要轻视四十天的时期，因为它包含了主行为的效法。在受难周之后，不要忽略在第四日和第六日禁食，同时将你的富裕分给穷人。若有人在主日或安息日禁食，除非是在逾越节安息日，他就是杀害基督的人。\par

% structure: p0106 source=h3 target=subsection reason=relative-heading-rank; base=h2

\subsection{第十四章　告别与警告}

你们的祈祷要延伸到安提约基雅的教会，因为我也作为囚犯正被带往罗马。我问候圣主教波利卡普；我问候圣主教味塔利乌，以及神圣的司铎团，和我的同仆执事们；愿我的灵魂取代他们。再次向主教和主内的司铎们告别。若有人同犹太人一同庆祝逾越节，或接受他们节日的象征物，他就是与杀害主及其宗徒的人有份。\par

% structure: p0108 source=h3 target=subsection reason=relative-heading-rank; base=h2

\subsection{第十五章　问候。结语}

斐罗执事和阿加托普斯执事问候你们。我问候贞女团体和寡妇序列；愿我因他们而喜乐！我问候主的子民，从最小的到最大的。我通过读者厄法尼乌将这封信送交你们，他是天主所尊重、极为忠实的人，我在勒基翁遇见他正要上船。请记念我的锁链\BibleRef{哥 4:18}，好使我在基督内得以成全。祝你们在肉身、灵魂和心神内安好，思想那完美的事，远离那些败坏真理之言、以主耶稣基督的恩宠在内里坚固的作恶者。\par

% structure: p0110 source=h2 target=section reason=relative-heading-rank; base=h2

\section{归化者玛利亚致依纳爵书}

% structure: p0111 source=h3 target=subsection reason=relative-heading-rank; base=h2

\subsection{问候}

玛利亚，耶稣基督的归化者，致安提约基雅宗徒教会最蒙福的主教依纳爵·德奥福鲁，在父天主和耶稣内蒙爱者：幸福与平安。我们都在祂内为您祈求喜乐与健康。\par

% structure: p0113 source=h3 target=subsection reason=relative-heading-rank; base=h2

\subsection{第一章　本书信的缘由}

由于基督已在我们中间奇妙地被显明为永生天主之子，并在末世藉童贞玛利亚、达味和亚巴郎的后裔成为人，正如先知团体先前关于祂、并通过祂所宣布的那样，因此我们恳切请求，藉您的智慧，我们的朋友马里斯，我们家乡内阿波利斯（靠近匝尔布斯）的主教，以及欧洛吉乌和司铎索贝鲁，能被派到我们这里，免得我们缺少主持圣言的人，正如梅瑟所说：\BibleQuote{愿主天主拣选一个人，引导这民众，上主的会众不要像没有牧人的羊群。}\BibleRef{户 27:16}\par

% structure: p0115 source=h3 target=subsection reason=relative-heading-rank; base=h2

\subsection{第二章　青年可与虔诚与明辨相结合}

至于我们所提到的那些年轻人，有福的人，你无需担忧。因为我愿你知道，他们在肉身方面是有智慧的，不为情欲所动，他们自身因着内在的德性，焕发着满头白发的全部荣耀，虽然他们作为年轻人刚被召叫担任司铎职。现在，请你藉天主藉基督赐给你的圣神，运用你的思想，你就会记起，撒慕尔还是幼童时，就被称为先见者，被列在先知的行列中；他指责年迈的厄里违犯诫命，因为厄里尊崇他那些愚蠢的儿子超过万物的创造者天主，并任由他们不受惩罚，即使他们嘲弄司铎职分，又粗暴对待你的百姓。\par

% structure: p0117 source=h3 target=subsection reason=relative-heading-rank; base=h2

\subsection{第三章 年轻献身的榜样}

此外，明智的达内尔年轻时曾审判几个刚愎的老人，显示他们是堕落之人，不配被视为长老，并且尽管按血统是犹太人，行为却是客纳罕人。耶肋米亚因自己年轻而推辞天主托付给他的先知职务时，被如此告诫：\BibleQuote{不要说：我太年轻，因为我派你到谁那里，你就应到谁那里去；我吩咐你说什么，你就应说什么。因为我和你在一起。}\BibleRef{耶 1:7}而明智的撒罗满年方十二岁时，就有智慧判断那两个女人关于孩子的疑难问题，当时众人不知孩子各属何人，以致全体人民都惊叹孩童身上的这种智慧，敬重他不仅是少年，更是成熟的人。他还解答了埃塞俄比亚女王的难题，那些难题如同尼罗河的支流般蕴藏益处，使得那位本身极有智慧的女人也无比惊异。\par

% structure: p0119 source=h3 target=subsection reason=relative-heading-rank; base=h2

\subsection{第四章 同一主题续}

还有天主所爱的约史雅，当他几乎还未清晰说话时，就揭露那些心怀邪灵的人言语虚假，欺骗百姓。他也揭露魔鬼的诡计，公开暴露那些不是神的偶像；尚在婴孩时就杀死了他们的司铎，推倒他们的祭台，用死尸玷污祭献之地，拆毁庙宇，砍伐木偶，打碎石柱，挖开恶人的坟墓，使恶者的痕迹不再存留。他如此热诚于敬神之事，证明自己是恶者的惩罚者，尽管他说话时仍像孩子般口齿不清。达味，既是先知又是君王，也是我们救主按肉身的根苗，年轻时就被撒慕尔傅油为王。\EditorialNote{撒慕尔纪上 16}因为他自己某处曾说：\BibleQuote{我在我兄弟中最小，在我父家是最年轻的。}\par

% structure: p0121 source=h3 target=subsection reason=relative-heading-rank; base=h2

\subsection{第五章 对依纳爵表达敬意}

但我若试图一一列举所有在年轻时蒙天主悦纳、被托付先知、司铎或君王职分的人，时间就不够了。所提的这些例子，足以唤起你们的记忆。但我恳求你们不要认为我倨傲或炫耀（如我这样写信）。因为我述说这些，并非为教导你，而只是提醒我在天主内的父亲。因我知道自己的身份，不敢与你们相比。我问候你神圣的圣职人员和在你牧养下管理的爱基督的人民。我们这里所有的信徒都问候你。有福的牧者，请祈祷，使我在天主前健康。\par

% structure: p0123 source=h2 target=section reason=relative-heading-rank; base=h2

\section{致纳波利（近匝尔布）的玛利亚书}

% structure: p0124 source=h3 target=subsection reason=relative-heading-rank; base=h2

\subsection{问候}

依纳爵，亦称号为德奥福禄，致藉至高天主圣父及为我等死的主耶稣基督的恩宠获得怜悯的、我的女儿玛利亚，至为忠信、堪配天主、心中怀着基督者，愿她在天主内满享幸福。\par

% structure: p0126 source=h3 target=subsection reason=relative-heading-rank; base=h2

\subsection{第一章 赞赏她的卓越与智慧}

的确，眼见胜过书写，因为眼见属于感官之一，它不仅通过传递友谊的证据尊荣接受者，而且通过自身所接受的，丰富了对更美好事物的渴望。但第二避难所，如常言所说，就是书写实践，我们藉着你的信德从如此遥远之处得到一个方便的港湾，因为通过信件我们了解到你内在的卓越。因为善良人的灵魂，最智慧的女子啊！如同最纯净的泉水；它们以美丽吸引过路人来喝，即使他们并不口渴。而你的智慧，如同一声命令，邀请我们分享那些在你灵魂中丰涌而出的神圣甘泉。\par

% structure: p0128 source=h3 target=subsection reason=relative-heading-rank; base=h2

\subsection{第二章 他自身的状况}

但是，有福的女子，我现在已不再是自己的主人，而是受制于他人，被许多敌手的变幻意志所驱策，一方面流徙，一方面被囚，一方面被捆锁。然而我并不在意这些。是的，藉着他们加于我的伤害，我愈发获得门徒的品格，好能到达耶稣基督。愿我享受为我预备的痛苦，因为\BibleQuote{现时的苦楚，与将来在我们身上要显示的光荣，是不能较量的。}\BibleRef{罗 8:18}\par

% structure: p0130 source=h3 target=subsection reason=relative-heading-rank; base=h2

\subsection{第三章 他已顺从其请求}

我欣然按你信中所求而行，对于你证实为有德行的那些人，我毫无疑虑。因为我相信，你是以属天主的判断、而非出于肉体的偏袒为他们作证的。你所引用的许多圣经章节令我极为喜悦，我读后对这事的任何疑念便不复存在。因为我并不认为，你已经向我提供了如此无可辩驳的证明、而我却只需用眼睛浏览一下即可。愿我成为你灵魂的替代，因为你爱耶稣，永生天主子。因此祂自己也对你说：\Quote{我爱那爱我的人；寻求我的，必得平安。}\par

% structure: p0132 source=h3 target=subsection reason=relative-heading-rank; base=h2

\subsection{第四章 称赞与劝勉}

如今我想到要提及的是，我在你于罗马与蒙福的林诺神父在一起时所听到的关于你的报告是真实的，而现在，配受真福的克莱孟——伯多禄和保禄的听众——已继承了他。到了现在，你的名声已增加百倍；愿你这妇人，还能更进一步。我极渴望到你那里去，好能与你同得安息；但\Quote{人的道路不由自己}。\BibleRef{耶 10:23}因为（我所处的）军营守卫阻碍了我的计划，不容我继续前行。而且，在我现在的状况中，我既不能做什么，也不能忍受什么。因此，我认为写信是朋友们相互鼓励的次要手段，我向你神圣的灵魂致意，恳求你进一步加强你的活力。因为我们目前的劳苦是微不足道的，而所期待的赏报是巨大的。\par

% structure: p0134 source=h3 target=subsection reason=relative-heading-rank; base=h2

\subsection{第五章 致辞与祝福}

要避开那些否认基督受难和祂降生成人的人；目前有许多人染此病症。但就其他方面而言，劝诫你是荒谬的，因为你在一切善工和善言上都是完美的，并且也能在基督内劝勉他人。请代向所有与你有相同心意、并持守他们在基督内之救恩的人问安。司铎和执事们，尤其是神圣的黑洛，向你问安。我的主人卡西安和他的妻子——我的姊妹——以及他们极亲爱的孩子们，也向你问安。愿主永远圣化你，使你在身体和灵魂的健康中获得享受，并愿我在基督内看见你获得冠冕！\par

% structure: p0136 source=h2 target=section reason=relative-heading-rank; base=h2

\section{致圣若望第一书}

依纳爵及与他同在的弟兄们，致书给神圣的司铎若望。\par

我们因你迟迟不来用讲论和安慰坚固我们而深感忧伤。若你的缺席再延长，将使我们的许多人失望。因此，请速速前来，因为我们相信这是有益的。此外，我们这里有许多妇女渴望见到耶稣的母亲玛利亚，她们日复一日想要离开我们到你那里去，好能与她相遇，触摸那喂养了主耶稣的乳房，并向她询问一些较为隐秘的事。还有你所爱的莎乐美（安娜之女），她曾在耶路撒冷与玛利亚同住了五个月，以及其他几位知名人士，都述说她满溢一切恩宠和美德，如同童贞女一般，在美德和恩宠上多结果实。据他们报告，她在迫害与患难中喜乐，在贫困与匮乏中毫无怨言，感激那些伤害她的人，在遭遇苦难时欢喜；她同情可怜和受苦的人，如同分担他们的痛苦，并且毫不迟疑地前去帮助他们。此外，她在信德的战斗中，与邪恶原则或行为的败坏冲突相对抗，荣光四射。她是我们新宗教和悔改的主母，是在信徒中作一切虔敬之事的谦卑婢女。她确实献身于谦卑的人，并且比献身者更谦卑地自卑自己，受到所有人的奇妙尊崇，同时也遭受经师和法利塞人的诽谤。除了这些，还有许多人告诉我们关于她的众多其他事情。然而，我们并不完全相信每件事的全部细节，也不向你提及这些。但根据可靠人士的告知，在耶稣的母亲玛利亚身上，有一种与人性结合的天使般的纯洁本性。这样的报告大大激动了我们的情绪，促使我们急切地渴望一睹这（若可如此说）属天的奇观和最神圣的奇迹。但愿你赶快满足我们的愿望。祝好。阿们。\par

% structure: p0139 source=h2 target=section reason=relative-heading-rank; base=h2

\section{致圣若望第二书}

他的朋友依纳爵致书给神圣的司铎若望。\par

若你允许，我愿上耶路撒冷去，看望那里的忠实圣徒，尤其是母亲玛利亚——他们报告说她为众人所钦佩和爱慕。因为，只要他是我们信仰和宗教的朋友，谁不愿看见并与那从自己胎中生育了真天主的她交谈呢？同样，我也愿见可敬的雅各伯——他别号义人——据说他在外貌、生活和行事方式上与基督耶稣极其相似，仿佛同一母腹的双生子。他们说，若我看见他，也如同看见了耶稣本人，满有祂身体的一切特征和容貌。此外，我也愿见其他男女圣徒。唉！为何我耽延？为何我被阻止？仁慈的老师，请吩咐我赶快（实现我的愿望），祝你平安。阿们。\par

% structure: p0142 source=h2 target=section reason=relative-heading-rank; base=h2

\section{依纳爵致童贞玛利亚书}

她的朋友依纳爵致书给怀基督的玛利亚。\par

你本应安慰并鼓励我——我是一名新信徒，也是你所爱的若望的门徒。因为我听到了关于你的儿子耶稣的奇妙事迹，这样的报告令我惊叹。但我全心渴望从你那里得知我所听到的事，因为你始终与祂亲密交好，并知晓祂的一切秘密。我也曾在另一时间写信给你，就同样的事询问你。祝你平安；愿与我同在的新信徒，因你、藉你、并在你内得到安慰。阿们。\par

% structure: p0145 source=h2 target=section reason=relative-heading-rank; base=h2

\section{童贞圣母对此信的回复}

基督耶稣的卑微婢女致书给依纳爵，她亲爱的同门。\par

从\Person{若望}那里所听见并学到的关于\Person{耶稣}的事，是真实的。你们要相信它们，持守它们，并坚守你们所接受的基督信仰的宣认，使你们的习惯和生活与你们的宣认相称。现在我将与\Person{若望}一同来探望你们和你们那里的众人。你们要坚守信德，\BibleRef{格前 16:13}，要作大丈夫；不要让迫害的凶猛动摇你们，却要使你们的精神刚强，并因天主你们的救主而喜乐。\BibleRef{路 1:47}。阿们。\par

\SourceRecord{Translated by Alexander Roberts and James Donaldson.；Ante-Nicene Fathers；Vol. 1.；Edited by Alexander Roberts, James Donaldson, and A. Cleveland Coxe.；Buffalo, NY: Christian Literature Publishing Co.,；1885.；<http://www.newadvent.org/fathers/0114.htm>.}
